\begin{examplebox}
\par\noindent\begin{minipage}{\linewidth}
  \textbf{\textcolor{CExample}{例 1（最基本的识别）：三边 $5,6,7$，没有高和角}}

  \textbf{题目：}三角形的三边长为 $5,6,7$，求面积。

\THeronFigure{5}{6}{7}{1}{0}{5}{6}{7}{例1：只给三边，不先求高或角，直接求阴影面积。}
\end{minipage}\par
  \textbf{方法选择：}三边齐全，没必要另设高。
  \[
   p=\frac{5+6+7}{2}=9,
  \]
  \[
   S=\sqrt{9(9-5)(9-6)(9-7)}
     =\sqrt{9\cdot4\cdot3\cdot2}=\boxed{6\sqrt6}.
  \]
  \textbf{检查：}$5+6>7$，三边确实能围成三角形；答案是正数。

  \medskip
\par\noindent\begin{minipage}{\linewidth}
  \textbf{\textcolor{CExample}{例 2（直接带出两个新量）：三边 $13,14,15$}}

  \textbf{题目：}三角形三边长分别为 $13,14,15$，求面积、内切圆半径以及最长边上的高。

\THeronFigure{13}{14}{15}{0.45}{4}{13}{14}{15}{例2：最长边 $AB=15$，它对应的高是 $CD=h$。}
\end{minipage}\par
  \textbf{第一步：半周长}
  \[
   p=\frac{13+14+15}{2}=21.
  \]
  \textbf{第二步：面积}
  \[
   S=\sqrt{21\cdot8\cdot7\cdot6}=\sqrt{7056}=\boxed{84}.
  \]
\THeronFigure{13}{14}{15}{0.43}{5}{13}{14}{15}{面积 $S$ 已求出；三块同高面积给出 $S=pr$，再求半径。}
  \textbf{第三步：内切圆半径}
  \[
   r=\frac Sp=\frac{84}{21}=\boxed4.
  \]
  \textbf{第四步：以最长边 $15$ 为底的高}
  \[
   84=\frac12\cdot15\cdot h
   \quad\Longrightarrow\quad h=\frac{168}{15}=\boxed{\frac{56}{5}}.
  \]
  \textbf{作用：}一旦面积已知，\textbf{内切圆半径和任意一边上的高}都可以直接求。

  \medskip
\par\noindent\begin{minipage}{\linewidth}
  \textbf{\textcolor{CExample}{例 3（反着用）：已知面积和周长，反求两条边}}

  \textbf{题目：}三角形面积为 $6$，周长为 $12$，其中一边为 $5$，求另外两边。

\THeronFigure{3}{4}{5}{1.35}{0}{x}{7-x}{5}{例3：面积为 $6$，另两边之和为 $7$；图示一种满足条件的形状。}
\end{minipage}\par
  \textbf{第一步：设另外两边为 $x,y$，把周长用上。}
  \[
    x+y=12-5=7,\quad y=7-x,\quad p=6.
  \]
  此时三边需满足三角形不等式：
  \[
    x+5>7-x\Longrightarrow x>1,\qquad
    (7-x)+5>x\Longrightarrow x<6.
  \]
  所以 $1<x<6$，并且 $y=7-x>0$。

  \textbf{第二步：海伦公式写成方程。}
  \[
    6^2=6(6-5)(6-x)(6-y).
  \]
  因为 $6-y=6-(7-x)=x-1$，故
  \[
    36=6(6-x)(x-1)
    \quad\Longrightarrow\quad(6-x)(x-1)=6.
  \]
  \textbf{第三步：解一元二次方程。}
  \[
    6x-6-x^2+x=6
    \Longleftrightarrow x^2-7x+12=0
    \Longleftrightarrow (x-3)(x-4)=0.
  \]
  得 $x=3$ 或 $x=4$；分别对应 $y=4$ 或 $y=3$。
  两组数均满足 $1<x<6$，且由 $3^2+4^2=5^2$ 可核对面积确为 $\tfrac12\cdot3\cdot4=6$。

  \textbf{答案：}另外两边为 $\boxed{3\text{ 和 }4}$。

  \medskip
\par\noindent\begin{minipage}{\linewidth}
  \textbf{\textcolor{CExample}{例 4（方法对照）：不是遇到三边就必须套海伦公式}}

  \textbf{题目：}等腰三角形三边为 $5,5,6$，求面积。

\begin{center}
\begin{tikzpicture}[scale=1.15,font=\small,line join=round]
\coordinate (A) at (0,0);\coordinate (B) at (6,0);\coordinate (C) at (3,4);\coordinate (D) at (3,0);
\fill[blue!12] (A)--(C)--(D)--cycle;
\fill[orange!18] (B)--(C)--(D)--cycle;
\draw[thick] (A)--node[above left=5pt] {$5$} (C)--node[above right=5pt] {$5$} (B)--cycle;
\draw[thick,CProof] (C)--node[midway,right=5pt] {$h$} (D);
\pic[draw,angle radius=3mm] {right angle=C--D--B};
\path (A)--node[midway,below=12pt] {$3$} (D)--node[midway,below=12pt] {$3$} (B);
\foreach \p/\where in {A/below left,B/below right,C/above,D/below}
 {\fill (\p) circle[radius=1.3pt];\node[\where=4pt] at (\p) {$\p$};}
\node[anchor=north,align=center,text width=11cm] at (3,-1.25)
 {等腰三角形的高平分底边：$AD=DB=3$，$h^2+3^2=5^2$。};
\end{tikzpicture}
\end{center}

\end{minipage}\par
  \textbf{方法一（初中方法，运算更直观）：}底边取 $6$，顶点向底边作高，底边分成 $3$ 和 $3$，由勾股定理
  \[
   h=\sqrt{5^2-3^2}=4,\qquad S=\frac12\cdot6\cdot4=\boxed{12}.
  \]
  \textbf{方法二（海伦公式核对）：}
  \[
   p=8,\quad S=\sqrt{8(8-5)(8-5)(8-6)}
    =\sqrt{8\cdot3\cdot3\cdot2}=\boxed{12}.
  \]
  \textbf{比较：}两种都对；等腰三角形作高能立刻一分为二时，直接底乘高可能更省事。\textbf{海伦公式是一条选择，不是强制步骤。}

  \medskip
\par\noindent\begin{minipage}{\linewidth}
  \textbf{\textcolor{CExample}{例 5（数字大也可能很快）：三边 $9,10,17$}}

  \textbf{题目：}求三角形面积和内切圆半径。

\THeronFigure{9}{10}{17}{0.45}{0}{9}{10}{17}{例5：三角形较扁，但 $9+10>17$；$p-17=1$ 仍为正数。}
\end{minipage}\par
  \textbf{第一步：}检查 $9+10>17$，成立。
  \[
   p=\frac{9+10+17}{2}=18,
   \quad (p-9,p-10,p-17)=(9,8,1).
  \]
  \textbf{第二步：}先利用整数乘积开方：
  \[
   S=\sqrt{18\cdot9\cdot8\cdot1}
    =\sqrt{1296}=\boxed{36},
   \qquad r=\frac{36}{18}=\boxed2.
  \]
  \textbf{观察：}虽然 $17$ 很长，但它离另外两边之和 $19$ 很近，$p-17=1$，实际计算反而简单。

  \medskip
\par\noindent\begin{minipage}{\linewidth}
  \textbf{\textcolor{CExample}{例 6（判断不要乱用）：已知两边和夹角}}

  \textbf{题目：}三角形有两边长 $8,10$，夹角为 $30^\circ$，求面积。

\begin{center}
\begin{tikzpicture}[scale=.7,font=\small,line join=round]
\coordinate (C) at (0,0);\coordinate (B) at (8,0);
\coordinate (A) at ({5*sqrt(3)},5);\coordinate (H) at ({5*sqrt(3)},0);
\fill[blue!10] (A)--(B)--(C)--cycle;
\draw[thick] (A)--(B)--node[midway,below=12pt] {$8$} (C)
 --node[midway,above left=5pt] {$10$} cycle;
\draw[gray,dashed] (B)--(H);
\draw[CProof,thick,dashed] (A)--node[midway,right=6pt] {$h=10\sin30^\circ$} (H);
\pic[draw,angle radius=3mm] {right angle=C--H--A};
\pic[draw,CProof,angle radius=10mm] {angle=B--C--A};
\node at (2,.55) {$30^\circ$};
\foreach \p/\where in {A/above,B/below left,C/below left,H/below right}
 {\fill (\p) circle[radius=1.4pt];\node[\where=4pt] at (\p) {$\p$};}
\node[anchor=north,align=center,text width=11cm] at (4.6,-1.7)
 {已知两边及其夹角：$S=\tfrac12\times8\times10\sin30^\circ$，无需先求第三边。};
\end{tikzpicture}
\end{center}

\end{minipage}\par
  因为夹角已经给出，直接用高中面积公式：
  \[
   S=\frac12\cdot8\cdot10\cdot\sin30^\circ
    =40\cdot\frac12=\boxed{20}.
  \]
  \textbf{为什么不先用海伦？}它需要三边，如果先用余弦定理求第三边、再套海伦，计算量明显增加。

  \medskip
  \textbf{\textcolor{CExample}{小练习（附答案）}}
  \begin{enumerate}[leftmargin=*,itemsep=0.2em]
   \item 三边 $4,5,6$，求面积。\quad\textbf{答案：}$\displaystyle \frac{15\sqrt7}{4}$（$p=15/2$）。
   \item 三边 $3,4,8$，能否用海伦公式求真实三角形面积？\quad\textbf{答案：不能，$3+4<8$，三边不能组成三角形。}
   \item 面积 $24$、半周长 $12$，求内切圆半径。\quad\textbf{答案：}$r=24/12=2$（已知确有这样的三角形时）。
  \end{enumerate}
\end{examplebox}
